\documentclass{amsart}
% For dealing with references we use the comment environment
\usepackage{verbatim}
\newenvironment{reference}{\comment}{\endcomment}
%\newenvironment{reference}{}{}
\newenvironment{slogan}{\comment}{\endcomment}
\newenvironment{history}{\comment}{\endcomment}

% For commutative diagrams we use Xy-pic
\usepackage[all]{xy}

% We use 2cell for 2-commutative diagrams.
\xyoption{2cell}
\UseAllTwocells

% We use multicol for the list of chapters between chapters
\usepackage{multicol}

% This is generally recommended for better output
\usepackage{lmodern}
\usepackage[T1]{fontenc}

% For cross-file-references

% Package for hypertext links:
\usepackage{hyperref}

% For any local file, say "hello.tex" you want to link to please

% Theorem environments.
%
\theoremstyle{plain}
\newtheorem{theorem}[subsection]{Theorem}
\newtheorem{proposition}[subsection]{Proposition}
\newtheorem{lemma}[subsection]{Lemma}

\theoremstyle{definition}
\newtheorem{definition}[subsection]{Definition}
\newtheorem{example}[subsection]{Example}
\newtheorem{exercise}[subsection]{Exercise}
\newtheorem{situation}[subsection]{Situation}

\theoremstyle{remark}
\newtheorem{remark}[subsection]{Remark}
\newtheorem{remarks}[subsection]{Remarks}

\numberwithin{equation}{subsection}

% Macros
%
\def\lim{\mathop{\mathrm{lim}}\nolimits}
\def\colim{\mathop{\mathrm{colim}}\nolimits}
\def\Spec{\mathop{\mathrm{Spec}}}
\def\Hom{\mathop{\mathrm{Hom}}\nolimits}
\def\Ext{\mathop{\mathrm{Ext}}\nolimits}
\def\SheafHom{\mathop{\mathcal{H}\!\mathit{om}}\nolimits}
\def\SheafExt{\mathop{\mathcal{E}\!\mathit{xt}}\nolimits}
\def\Sch{\mathit{Sch}}
\def\Mor{\mathop{\mathrm{Mor}}\nolimits}
\def\Ob{\mathop{\mathrm{Ob}}\nolimits}
\def\Sh{\mathop{\mathit{Sh}}\nolimits}
\def\NL{\mathop{N\!L}\nolimits}
\def\CH{\mathop{\mathrm{CH}}\nolimits}
\def\proetale{{pro\text{-}\acute{e}tale}}
\def\etale{{\acute{e}tale}}
\def\QCoh{\mathit{QCoh}}
\def\Ker{\mathop{\mathrm{Ker}}}
\def\Im{\mathop{\mathrm{Im}}}
\def\Coker{\mathop{\mathrm{Coker}}}
\def\Coim{\mathop{\mathrm{Coim}}}

% Boxtimes
%
\DeclareMathSymbol{\boxtimes}{\mathbin}{AMSa}{"02}

%
% Macros for moduli stacks/spaces
%
\def\QCohstack{\mathcal{QC}\!\mathit{oh}}
\def\Cohstack{\mathcal{C}\!\mathit{oh}}
\def\Spacesstack{\mathcal{S}\!\mathit{paces}}
\def\Quotfunctor{\mathrm{Quot}}
\def\Hilbfunctor{\mathrm{Hilb}}
\def\Curvesstack{\mathcal{C}\!\mathit{urves}}
\def\Polarizedstack{\mathcal{P}\!\mathit{olarized}}
\def\Complexesstack{\mathcal{C}\!\mathit{omplexes}}
% \Pic is the operator that assigns to X its picard group, usage \Pic(X)
% \Picardstack_{X/B} denotes the Picard stack of X over B
% \Picardfunctor_{X/B} denotes the Picard functor of X over B
\def\Pic{\mathop{\mathrm{Pic}}\nolimits}
\def\Picardstack{\mathcal{P}\!\mathit{ic}}
\def\Picardfunctor{\mathrm{Pic}}
\def\Deformationcategory{\mathcal{D}\!\mathit{ef}}

\setlength{\emergencystretch}{3em}
\begin{document}
\title[Stacks Project, Tag 03RD]{279 --- Tsen-Lang theorem for function fields of r-dimensional varieties over algebraically closed fields}
\maketitle
\noindent Copyright (C) 2005--2025 Johan de Jong. Stacks Project authors. Source: \href{https://stacks.math.columbia.edu/tag/03RD}{Tag 03RD}.

\noindent Editorial difficulty note (not author text). Difficulty H2: The central task is to construct a nonzero common solution by replacing rational unknowns with global sections of increasing ample powers. Asymptotic Riemann-Roch must compare the growth of the coefficient equations and unknown sections with the sharp inequality n>d\textasciicircum{}r, reducing the problem to homogeneous intersection over an algebraically closed field..

\noindent Editorial provenance note (not author text). History: excerpt from revision \texttt{a0444\allowbreak 6e57e\allowbreak c1fbc\allowbreak 252a8\allowbreak 71afc\allowbreak ec775\allowbreak 2fb28\allowbreak 07b14}, etale-cohomology.tex, lines 10861--10891. Reference presentation was adapted; original-author context and core support blocks were retained or added. The mathematical wording is unchanged. Licensed under GNU FDL 1.2 or later; no invariant sections or cover texts. See ../licenses/COPYING. Context blocks reproduce author definitions and supporting results; they are not additional counted problems.

\begin{theorem}[Tsen's theorem]
\label{theorem-tsen}
The function field of a variety of dimension $r$ over an algebraically closed
field $k$ is $C_r$.
\end{theorem}

\begin{proof}
For projective space one can show directly that the field
$k(x_1, \ldots, x_r)$ is $C_r$ (exercise).

\medskip\noindent
General case. Without loss of generality, we may assume $X$ to be projective.
Let $f \in k(X)[T_1, \ldots, T_n]_d$ with $0 < d^r < n$.
Say the coefficients of $f$ are in $\Gamma(X, \mathcal{O}_X(H))$
for some ample $H \subset X$. Let
$\mathbf{\alpha} = (\alpha_1, \ldots, \alpha_n)$ with $\alpha_i \in \Gamma(X,
\mathcal{O}_X(eH))$. Then $f(\mathbf{\alpha}) \in \Gamma(X,
\mathcal{O}_X((de + 1)H))$. Consider the system of equations $f(\mathbf{\alpha})
=0$. Then by asymptotic Riemann-Roch
(Varieties, Proposition \href{https://stacks.math.columbia.edu/tag/0BJ8}{proposition asymptotic riemann roch (Tag 0BJ8)})
there exists a $c > 0$ such that
\begin{itemize}
\item the number of variables is
$n\dim_k \Gamma(X, \mathcal{O}_X(eH)) \sim n e^r c$, and
\item the number of equations is
$\dim_k \Gamma(X, \mathcal{O}_X((de + 1)H)) \sim (de + 1)^r c$.
\end{itemize}
Since $n > d^r$, there are more variables than equations. The equations are
homogeneous hence there is a solution by
Lemma \href{https://stacks.math.columbia.edu/tag/03RB}{lemma algebraically closed find solutions (Tag 03RB)}.
\end{proof}
\end{document}
